% Epilogue of Optimizers in the Wild.
\specialchapter{Epilogue}{What theory owes practice}{30}\label{ch:epilogue}
\begin{refsection}[chapters/16-epilogue.bib]
% Tables outside the numbered chapters are numbered in one plain sequence: the prologue has
% Tables 1 and 2, so the ledger is Table 3. The redefinition is local to this environment.
\renewcommand{\thetable}{\arabic{table}}\setcounter{table}{2}

The prologue listed the decisions that a practitioner makes before a training run and asked what theory can say about each. This epilogue keeps the account. It records, practice by practice, what the chapters proved and in which idealization, what they predicted, and what they left unexplained; it says what the idealizations miss when they are taken together; and it lists the questions that the book raises and cannot answer.

\addsec{A ledger}

\Cref{tab:epilogue:ledger} has one row for each practice. The verdict is one of three. A practice is \emph{explained} when a theorem in a faithful idealization accounts for it and the prediction survives a comparison with simulation or with the evidence from training runs. It is \emph{partly explained} when a mechanism has been identified and proved in an idealization, but either the idealization is far from training or the prediction is only qualitative. It is \emph{unexplained} when the book has nothing to say beyond a description, or when the practice is folklore. The last column says what the verdict rests on, which is where a reader who doubts it should look.

\begin{table}[p]
\begin{fullwidth}
\centering\footnotesize
\renewcommand{\arraystretch}{1.14}
\begin{tabular}{@{}>{\raggedright\arraybackslash}p{31mm}r>{\raggedright\arraybackslash}p{17mm}>{\raggedright\arraybackslash}p{104mm}@{}}
\toprule
\textsc{practice} & \textsc{ch.} & \textsc{verdict} & \textsc{what the verdict rests on}\\
\midrule
Normalized updates & \ref{ch:steepest} & explained & Every update in use is a steepest-descent step in some norm; rates depend on the norm only through a smoothness constant and a dual norm\\
Sign-like steps (Adam, Lion) & \ref{ch:sign} & explained & Memoryless Adam and Lion take sign steps; with memory Adam's step is at most $C(\beta_1,\beta_2)$ learning rates, $7.27$ at $(0.9,0.999)$ and $1.16$ at $(0.9,0.95)$\\
When $\ell_\infty$ steps help & \ref{ch:sign} & partly & A gradient denser than the effective rank of a curvature that is nearly diagonal in the parameters' coordinates; the comparison is between guarantees\\
Muon's orthogonalization & \ref{ch:spectral} & explained & The steepest step in the spectral norm is the polar factor; Muon's quintic does not converge but holds all but the smallest singular values in $[0.681,1.203]$\\
When Muon helps & \ref{ch:spectral} & partly & Nuclear rank of the gradient against stable rank of the inputs, for one-step guarantees; many steps and the shrinking gains at scale are not explained\\
Width rules: $\mu$P, Muon's shape factor & \ref{ch:modular} & partly & Width exponents follow from an aligned step in the RMS$\to$RMS norm and match measured slopes, but alignment fades; base rates are not derived\\
Momentum & \ref{ch:momentum} & partly & Stability regions and stationary variances are exact on quadratics; normalized steps need memory when the noise is skewed; why $\beta=0.9$ is not explained\\
Iterate and weight averages & \ref{ch:momentum} & partly & For convex losses a uniform average at a constant step matches a decayed step; on quadratics a weight average cools only the stiff directions (derived)\\
Adam's $\beta_2$ & \ref{ch:adaptive} & partly & Rare gradients need a long memory $1/(1-\beta_2)$, which costs an excess move of about $2\ln2/(1-\beta_2)$ learning rates after a quiet spell; the values are not derived\\
Shampoo, K-FAC, SOAP & \ref{ch:kronecker} & partly & The Kronecker bound and Shampoo's regret are proved, but the regret never beats an adaptive gradient step; the best power depends on what the factors estimate; the wins over AdamW are not derived\\
Training at $2/\eta$ & \ref{ch:eos} & partly & Thresholds for gradient descent, momentum, and Adam with a frozen preconditioner are exact on quadratics; self-stabilization is proved in a cubic model, in continuous time; the rise of sharpness is an input\\
Linear warmup & \ref{ch:warmup} & partly & If curvature is bounded by the loss, then a warmup pays the initial curvature once and a constant step at every scale, and the best stable rise is hyperbolic, not linear; the bound is measured, not derived, and the length is not explained\\
Gradient clipping & \ref{ch:spikes} & partly & Under $(L_0,L_1)$-smoothness, measured but not derived, clipping avoids paying for the largest gradient, which no fixed step can; the threshold $1.0$ is folklore\\
z-loss and QK-norm & \ref{ch:spikes} & partly & Each removes a mechanism of instability (a flat direction of the softmax; attention logits quadratic in the weights), though a z-loss of $10^{-4}$ is too weak to pin a unigram model's logits; real spikes have interacting causes\\
Batch-size rules & \ref{ch:noise} & partly & In the noisy quadratic model the best step grows linearly with the batch up to a critical batch size of order $\tr\Sigma/(S\eps)$, which preconditioning raises; Adam's square-root rule is derived, not proved; curvature and noise are frozen\\
Cosine and WSD & \ref{ch:schedules} & partly & Convex last-iterate bounds for every schedule, and an exact river valley with a noise floor and a cooldown law; with tuned peaks the valley ranks cosine well ahead of WSD, while in language models the two match\\
Decoupled weight decay & \ref{ch:decay} & explained & Steepest descent with decoupled decay is Frank--Wolfe on a norm ball; under SGD scale-invariant layers settle at a rotation of $\arccos(1-\eta\lambda)$ per step, as vision runs confirm; momentum shrinks the angle\\
The value $\lambda=0.1$ & \ref{ch:decay} & unexplained & Only the product $\eta\lambda$ and the timescale $1/(\eta\lambda)$ matter; neither value is derived\\
Adam over SGD for transformers & \ref{ch:transformers} & partly & Three solvable mechanisms: heavy-tailed classes slow gradient descent to a logarithmic decay; blocks of curvature at different scales; and a gap that grows with the batch size; their interaction is open\\
Adaptive methods and generalization & \ref{ch:bias} & partly & Steepest descent ends at the maximum margin of its own norm on separable data, a prior that can favor either optimizer; the evidence for SGD is mixed\\
\bottomrule
\end{tabular}
\caption{The practices of the book, with the chapter that treats each, its verdict, and what the verdict rests on.}
\label{tab:epilogue:ledger}
\end{fullwidth}
\end{table}

Three things stand out. First, the update rules themselves are well explained: the first line of the book, that optimizers are steepest descent in norms with preconditioning and moving averages, accounts for the form of every rule in the field guide, with decoupled weight decay read as a Frank--Wolfe step, and it predicts how learning rates should scale with width while updates stay aligned with the features. Second, almost every practice of the second line is partly explained, and in the same way: a mechanism is proved in an idealization, and the idealization contains one or two parameters, a curvature bound, a noise level, a rate of sharpening, that are measured rather than derived. Third, the numbers that practitioners set in the recipes this book quotes, $\beta_1=0.9$, $\beta_2=0.95$, weight decay $0.1$, a clipping threshold of $1.0$, a warmup of two thousand steps, a cooldown of a fifth of training, are unexplained. The theory says which combinations matter and how they should scale; it does not say why these values.

\addsec{What the idealizations miss together}

Each chapter listed what its own idealization misses. Read together, the lists repeat five omissions.

\paragraph{The landscape is fixed.} Almost every idealization in the book runs an optimizer on a landscape given in advance: a quadratic, a convex function, a cubic model of the edge of stability, a valley with a river. Along the trajectory of a network the local landscape changes: the sharpness rises, features form, norms grow and settle, and the gradient noise changes with them. Where the book needed such a change, it put it in by hand, as the rate of progressive sharpening in \cref{ch:eos} or the loss-dependent curvature in \cref{ch:warmup}.

\paragraph{One mechanism at a time.} Warmup is explained in \cref{ch:warmup} by curvature that falls with the loss, by the edge of stability, and by the statistics of Adam's first steps, and in \cref{ch:decay} by the size of early angular updates; the advantage of Adam on transformers by three mechanisms in \cref{ch:transformers}. In a training run they all operate at once, and nothing in the book says how they combine.

\paragraph{Guarantees, not performance.} Several comparisons in the book are comparisons of worst-case guarantees: which norm gives the better descent bound (\cref{ch:steepest,ch:sign,ch:spectral}), which schedule gives the better last-iterate bound (\cref{ch:schedules}). Where the chapters measured how far a guarantee is from what the optimizer does, the answer ranged from nothing, for sign descent on a diagonal quadratic, to more than a factor of ten, for AdaGrad's regret and for normalized momentum (\cref{ch:steepest,ch:adaptive,ch:momentum}). A guarantee for one step can also misrank two methods over a run: on generic targets the first spectral step beats the first gradient step, yet spectral descent can need more than twice as many steps to reach the same loss (\cref{ch:spectral}).

\paragraph{Simple noise.} The noise in most chapters has bounded variance, often Gaussian, often proportional to the curvature. Heavy tails get a theorem only in \cref{ch:spikes}, where clipping restores a rate, and an exercise in \cref{ch:momentum}. The structure of the data enters through rare classes (\cref{ch:sign,ch:spikes}), sparse features (\cref{ch:adaptive}), and the Zipf law of token frequencies (\cref{ch:transformers}), and each time it changed the answer: a sign method without memory ignores a rare class, and gradient descent slows to a logarithmic decay on Zipfian data.

\paragraph{No scale.} A careful benchmark found that the speedups of newer optimizers over a well-tuned AdamW shrink as models grow, for the fastest of them from about $1.4\times$ at $0.1$ billion parameters to about $1.1\times$ at $1.2$ billion \parencite{WenEtAl2026Fantastic}. No idealization in this book has a size parameter that reproduces this. The nearest candidate is the noisy quadratic model of \cref{ch:noise,ch:transformers}: there the advantage of a preconditioner falls from its condition-number factor at large batch sizes to a noise-weighted constant at small ones, and the critical batch size grows as the target loss falls, so that a run with a fixed batch size moves toward the small-batch end as it aims lower. The book does not test whether this is why the gains shrink.

\addsec{What is open}

The chapters end with research exercises. Here are the questions that recur across them, in roughly the order of the book.

\begin{enumerate}
\item \emph{The right norm.} The norm of a network is assembled from norms on its layers, but the norms for embeddings, for attention, for normalization gains, and for the output layer are chosen, not derived (\cref{ch:spectral,ch:modular}). Is there a principle, perhaps from the statistics of activations and gradients, that selects them?
\item \emph{The averaging constants.} Why do $\beta_1=0.9$ and $\beta_2\in[0.95,0.999]$ work so widely (\cref{ch:momentum,ch:adaptive})? The theory explains which timescales these constants set, not why those timescales are right.
\item \emph{Progressive sharpening.} Full-batch gradient descent on a network raises the sharpness until the step is barely stable (\cref{ch:eos}). Why does the sharpness rise in the first place? A model in which progressive sharpening is an output rather than an input would bear on most of Parts III and IV.
\item \emph{Long warmups.} The curvature model of \cref{ch:warmup} sets the length of a warmup by the initial excess curvature, about $K_1\Delta_0/\mu$ steps, which is unknown for a network, and Adam's statistics suggest $2/(1-\beta_2)$ steps, forty at $\beta_2=0.95$, while LLaMA, with that $\beta_2$, warms up for $2000$. What sets the length?
\item \emph{Cooldowns.} The river valley of \cref{ch:schedules} predicts the noise floor that a cooldown removes and how the best cooldown fraction scales with the length of the run, but through an exponent and a spread of the noise that must be measured; with tuned peaks it ranks cosine well ahead of warmup--stable--decay, while in language models the two match. What limits the peak, and what sets a cooldown near a fifth of the run?
\item \emph{Shrinking gains.} Why do the gains of Muon, SOAP, and other matrix methods over AdamW shrink with model size (\cref{ch:spectral,ch:kronecker})? An unrefereed preprint reports that Hyperball, which fixes the norms of weights and updates, gives Muon token-equivalent speedups of $20$ to $30\%$ over weight-decay baselines up to $1.2$ billion parameters \parencite{WenEtAl2026Hyperball}; whether such gains survive at larger scale, and whether the effective learning rate that weight decay sets explains the difference (\cref{ch:decay}), are open.
\item \emph{Hyperparameter scaling laws.} How should the learning rate, the batch size, and the weight decay scale with model size and data size together? The book gives the exponent of width (\cref{ch:modular}), the critical batch size of a quadratic (\cref{ch:noise}), and the timescale of weight decay (\cref{ch:decay}); a joint theory is open.
\item \emph{Spikes.} A solvable model of a loss spike in a transformer, with the optimizer state, the data, and the architecture all playing their parts, is open (\cref{ch:spikes}).
\item \emph{Beyond pretraining.} The optimizers of this book are tuned for pretraining. Fine-tuning and reinforcement learning are not treated, mixtures of experts appear only through experts that a router stops choosing (\cref{ch:spikes}), and low precision only through rounding in bfloat16 (\cref{ch:spectral,ch:spikes}); each raises questions that the book does not answer.
\end{enumerate}

\addsec{How to turn a practice into a theorem}

The chapters followed one method, and it may be useful to state it. Begin with the practice as practitioners state it, with its numbers and its sources, and with the evidence for and against it; the first section of every chapter is an attempt at this, and the attempt often showed that the practice was folklore. Find the smallest idealization in which the practice has a reason: one coordinate for sign descent, one matrix for Muon, a cubic for the edge of stability, a valley with a river for schedules, a sphere for weight decay. Prove what the idealization implies, completely, so that its consequences are not in doubt. Extract a number from the proof: a threshold, an exponent, a rate, a timescale, an equilibrium. Compare the number with a simulation, and then with training runs, and report the disagreement. Finally, list what the idealization threw away, and look in that list for the next idealization.

The method has a limit that the ledger makes plain: it explains mechanisms one at a time, in models that are much smaller than the networks they describe. Whether the mechanisms add up to an account of large-scale training is the question this book leaves open.

\notesandrefs

The benchmark of \textcite{WenEtAl2026Fantastic} compares ten optimizers on language models of $0.1$ to $1.2$ billion parameters, trained on one to eight times the Chinchilla-optimal amount of data, with careful tuning of each, and finds that the speedups over AdamW shrink with scale; its follow-up \parencite{WenEtAl2026Hyperball}, an unrefereed preprint, proposes fixing the norms of weights and updates directly. Several practices of the ledger, among them warmup, clipping, the batch size, and the schedule, are discussed from the practitioner's side in the \emph{Deep Learning Tuning Playbook} \parencite{GodboleEtAl2023}. The individual chapters give the sources for each entry of the ledger.

\printchapterbib
\end{refsection}
